\chapter*{PHỤ LỤC}
\addcontentsline{toc}{chapter}{PHỤ LỤC}
\markboth{PHỤ LỤC}{}

Phụ lục liệt kê tài nguyên và chi tiết kỹ thuật phục vụ việc tái lập kết quả đề tài.

\section*{A. Tài nguyên công khai}
\addcontentsline{toc}{section}{A. Tài nguyên công khai}

\begin{table}[H]
\centering
\caption{Tài nguyên sử dụng trong đề tài}
\label{tab:tainguyen}
\small
\begin{tabular}{p{4.0cm} p{9.5cm}}
\toprule
\textbf{Tài nguyên} & \textbf{Địa chỉ / định danh} \\
\midrule
Bộ dữ liệu & \url{https://www.kaggle.com/datasets/pkdarabi/cardetection} \\
Trọng số mô hình thầy & \texttt{weights/yolo26/teacher\_yolo26s.pt} \\
Trọng số mô hình trò & \texttt{weights/yolo26/student\_\{baseline,kd\}\_yolo26n.pt} \\
Bản lượng tử hoá & \texttt{weights/yolo26/student\_kd\_int8.onnx} (616 nút lượng tử hoá) \\
Thư viện phát hiện & Ultralytics YOLO \cite{jocher2023ultralytics} \\
Bộ tính độ đo & pycocotools trên định dạng COCO \cite{lin2014coco} \\
Nền tảng học sâu & PyTorch \cite{paszke2019pytorch} \\
Môi trường huấn luyện & Kaggle Notebook, GPU Tesla T4 \\
Sổ tay thực nghiệm & \texttt{notebooks/giaidoan2.ipynb} \\
\bottomrule
\end{tabular}
\end{table}

\section*{B. Danh sách lớp}
\addcontentsline{toc}{section}{B. Danh sách lớp}

Mười lăm lớp được giữ nguyên theo bộ dữ liệu gốc, không gộp và không ánh xạ lại:

\begin{table}[H]
\centering
\caption{Danh sách 15 lớp và mã tương ứng}
\label{tab:danhsachlop}
\small
\begin{tabular}{clcl}
\toprule
\textbf{Mã} & \textbf{Tên lớp} & \textbf{Mã} & \textbf{Tên lớp} \\
\midrule
0 & Green Light      & 8  & Speed Limit 40 \\
1 & Red Light        & 9  & Speed Limit 50 \\
2 & Speed Limit 10   & 10 & Speed Limit 60 \\
3 & Speed Limit 100  & 11 & Speed Limit 70 \\
4 & Speed Limit 110  & 12 & Speed Limit 80 \\
5 & Speed Limit 120  & 13 & Speed Limit 90 \\
6 & Speed Limit 20   & 14 & Stop \\
7 & Speed Limit 30   &    & \\
\bottomrule
\end{tabular}
\end{table}

\section*{C. Cách chạy lại quy trình}
\addcontentsline{toc}{section}{C. Cách chạy lại quy trình}

Mỗi bước của quy trình là một mô-đun độc lập, gọi được trực tiếp từ dòng lệnh:

\begin{lstlisting}[language=bash, caption={Các lệnh chạy lại quy trình}, label={lst:cmds}]
# 1. Kiem tra du lieu, dong bo cau hinh, kiem toan ro ri
python -m src.data.inspect_dataset --write-configs
python -m src.data.validate_labels
python -m src.data.audit_leakage

# 2. Phan tich kham pha du lieu
python -m src.eda.class_distribution
python -m src.eda.bbox_statistics
python -m src.eda.image_statistics
python -m src.eda.heatmap_analysis

# 3. Huan luyen ba giai doan  (can GPU, ~3 gio tren Tesla T4)
python -m src.compression.distill --stage teacher
python -m src.compression.distill --stage baseline
python -m src.compression.distill --stage kd

# 4. Xuat ONNX va luong tu hoa INT8
python -m src.compression.quantize

# 5. Danh gia bang pycocotools, co AP tach theo kich thuoc
python -m src.data.convert_to_coco
python -m src.evaluation.coco_eval                 # toan tap kiem tra
python -m src.evaluation.coco_eval --clean-only    # chi 573 anh khong ro ri
python -m src.evaluation.speed

# 6. Kiem chung tren video duong that
python scripts/eval_videos.py --stride 10

# 7. Bo kiem thu tu dong  (chi can CPU)
pytest
\end{lstlisting}

\section*{D. Bảng ánh xạ hình vẽ và nguồn dữ liệu}
\addcontentsline{toc}{section}{D. Bảng ánh xạ hình vẽ và nguồn dữ liệu}

Mọi hình và bảng số liệu trong báo cáo đều được sinh ra từ dữ liệu thật, không có số liệu minh
hoạ. Bảng~\ref{tab:anhxa} chỉ rõ nguồn của từng nhóm.

\begin{table}[H]
\centering
\caption{Ánh xạ giữa nội dung báo cáo và tệp kết quả sinh ra bởi quy trình}
\label{tab:anhxa}
\small
\begin{tabular}{p{5.6cm} p{7.9cm}}
\toprule
\textbf{Nội dung trong báo cáo} & \textbf{Nguồn} \\
\midrule
Bảng thống kê tổng quan     & \texttt{results/tables/dataset\_summary.csv} \\
Bảng và hình phân bố lớp    & \texttt{results/tables/class\_distribution.csv} \\
Bảng phân nhóm kích thước   & \texttt{results/tables/bbox\_size\_categories.csv} \\
Thống kê ảnh                & \texttt{results/tables/image\_statistics.csv} \\
Báo cáo chất lượng nhãn     & \texttt{results/tables/data\_quality\_report.csv} \\
Ma trận nhầm lẫn            & Thư mục chạy của Ultralytics \\
Bảng thành phần bộ dữ liệu  & Phân tích bổ sung theo quy ước đặt tên tệp \\
Toàn bộ bảng kết quả, $AP$ theo kích thước, độ trễ, kiểm toán rò rỉ & \texttt{notebooks/giaidoan2.ipynb} \\
Đo lại tại chỗ, tập đầy đủ và tập sạch & \texttt{src/evaluation/coco\_eval.py} $\rightarrow$ \texttt{results/tables/coco\_eval\{,\_clean\}.csv} \\
Độ trễ đo lại tại chỗ & \texttt{src/evaluation/speed.py} $\rightarrow$ \texttt{results/metrics/speed.json} \\
Kiểm chứng trên video đường thật & \texttt{scripts/eval\_videos.py} $\rightarrow$ \texttt{results/video\_eval\_yolo26/} \\
Kiểm toán rò rỉ & \texttt{src/data/audit\_leakage.py} $\rightarrow$ \texttt{results/metrics/data\_leakage.json} \\
\bottomrule
\end{tabular}
\end{table}

\section*{E. Ghi chú về tính tái lập}
\addcontentsline{toc}{section}{E. Ghi chú về tính tái lập}

Hạt ngẫu nhiên mặc định là 42, được đặt cho biến môi trường băm của Python, thư viện ngẫu nhiên
chuẩn, NumPy, PyTorch và PyTorch CUDA. Chế độ tất định của cuDNN được bật và chế độ tự dò thuật
toán nhanh nhất được tắt --- điều này làm chậm huấn luyện một chút nhưng là điều kiện cần để kết
quả lặp lại được giữa các lần chạy.

Cần lưu ý rằng ngay cả khi đã gieo hạt đầy đủ, kết quả vẫn có thể sai khác nhỏ giữa các phiên bản
thư viện, các kiến trúc GPU khác nhau và giữa chế độ CPU với GPU. Các con số trong báo cáo được đo
trên môi trường Kaggle Notebook với GPU Tesla T4 tại thời điểm thực hiện đề tài.
